% GENERATED FILE. Edit canonical lessons, then run npm run generate:books.
% canonical-source-hash: fnv1a64:4bad10ec81f06f52
% canonical-lessons: ZH-C188-youqian, ZH-W188-qian-money, ZH-C188-huaqian, ZH-C188-lingqian, ZH-C188-shengqian, ZH-C188-gou

\chapter{Money to Spend}
\label{ch:zh-money-to-spend}
\hlchaptermodality{188}

\begin{chapteropening}
\textbf{By the end of this chapter:} \emph{I can say someone is well off, spend money or spend less, ask for change and say it is enough.}

Complete the last lesson of chapter 188: I can say someone is well off, spend money or spend less, ask for change and say it is enough.
\end{chapteropening}

\section[\texorpdfstring{\zh{有钱} (yǒuqián)}{有钱 (yǒuqián)}]{\zh{有钱} (yǒuqián) — rich; well-off}

\label{lesson:ZH-C188-youqian}



\begin{quote}
Before the new one: say the Mandarin for at that time, then the Mandarin for in the end.
\end{quote}

\subsection*{You'll want to know: \zh{有钱}}
\textbf{\zh{有钱}} — \emph{yǒuqián} — \textquotedblleft{}rich; well-off\textquotedblright{}.

\begin{quote}
\textbf{\zh{他很有钱}} \emph{tā hěn yǒuqián} — He's very well off.
\end{quote}

The new character, in pieces: \zh{钱} — ten strokes in two pieces: on the left the metal piece, five strokes; on the right two short bars, a long slanting hook, a stroke falling to the left across it, and a dot. The next lesson puts it in your hand.

\begin{cousinweb}
\textbf{\zh{有钱}} (\emph{yǒuqián}) is built from \zh{有} 'to have' and \zh{钱} 'money', both inherited from Middle Chinese, \emph{hjuwX} and \emph{dzjen} (Guangyun, in Baxter's transcription). The dictionary explains the second, after an older commentary, as money named after a spring, because it flows everywhere.
\end{cousinweb}

\begin{grammarlens}[title={What you've built}]
One more word for asking for what I want.
\end{grammarlens}

\subsection*{Your turn}
\begin{itemize}
\raggedright
  \item \emph{Say it:} \emph{yǒuqián}
  \item \emph{Say it:} \emph{yǒuqián}, once more
  \item \emph{Say it:} \emph{tā hěn yǒuqián} — He's very well off.
  \item \emph{Recall:} say the Mandarin for at that time, then the Mandarin for in the end, then say \emph{yǒuqián} again
\end{itemize}

\begin{tcolorbox}[breakable,colback=teal!4,colframe=teal!35!black,title={Before you move on}]
What does \zh{有钱} mean? (\textquotedblleft{}Rich; well-off\textquotedblright{}.) Say it once more.
\end{tcolorbox}
\section[\texorpdfstring{\zh{钱} (qián)}{钱 (qián)}]{\zh{钱} — the metal piece, then two bars and a hook}

\label{lesson:ZH-W188-qian-money}



\begin{quote}
The word first, then the character inside it.

\begin{itemize}
\raggedright
  \item \emph{Recall:} say \emph{yǒuqián} — \textbf{R1}, one lesson back, before it decays
\end{itemize}

\begin{quote}
\zh{有} \zh{多}
\end{quote}

\emph{Read it:} each one aloud, with nothing said first — \emph{yǒu}, \emph{duō}
\end{quote}

\subsection*{Script — the metal piece, then two bars and a hook}
\hlblockfigure{\detokenize{figures/ZH-W188-qian-money-filmstrip.pdf}}{How \zh{钱} is written, stroke by stroke}

\begin{quote}
\zh{钱}
\end{quote}

\textbf{Two pieces, left before right.} Count the pieces, not the strokes.

\textbf{The left piece, the metal piece, five strokes.} A long stroke falling to the left, two short bars, a short bar rising to the right, then a vertical that turns up to the right at its foot.

\textbf{The right piece, five strokes.} Two short bars, one over the other. A long slanting hook from the top down to the lower right, flicking up at its foot. A stroke falling to the left across it. Then a dot at the top right.

5 + 5 = ten.

\subsection*{Your turn}
\begin{itemize}
\raggedright
  \item \emph{Write it:} \zh{钱}, all ten, with the model in front of you
  \item \emph{Cover:} write it again from memory, and count the pieces, not the strokes
\end{itemize}

\begin{tcolorbox}[breakable,colback=teal!4,colframe=teal!35!black,title={Before you move on}]
How many strokes in \zh{钱}? (\textbf{Ten}.) What is the left piece? (\textbf{The metal piece, five strokes}.) What comes last? (\textbf{The dot at the top right}.)

Source: \href{https://raw.githubusercontent.com/chanind/hanzi-writer-data/68d10a4b21150cae5e1ebbd223eed289cf32d90c/data/\%E9\%92\%B1.json}{Hanzi Writer Data \zh{钱}, pinned stroke paths and medians}.
\end{tcolorbox}
\section[huāqián]{huāqián — to spend money}

\label{lesson:ZH-C188-huaqian}



\begin{quote}
Before the new one: say the Mandarin for in the end, then the Mandarin for well-off.
\end{quote}

\subsection*{You'll want to know: huāqián}
\emph{huāqián} — \textquotedblleft{}to spend money\textquotedblright{}.

\begin{quote}
\emph{wǒ mèimei hěn huì huāqián} — My little sister really knows how to spend money.
\end{quote}

Its characters come later; for now it is a word for the ear.

\begin{cousinweb}
\emph{huāqián} is built from \emph{huā}, the word for a flower, here 'to spend', and \zh{钱} 'money'. The second is inherited from Middle Chinese \emph{dzjen} (Guangyun, in Baxter's transcription); the dictionary lists the first only as a flower, so how it came to mean spending is not settled.
\end{cousinweb}

\begin{grammarlens}[title={What you've built}]
One more word for asking for what I want.
\end{grammarlens}

\subsection*{Your turn}
\begin{itemize}
\raggedright
  \item \emph{Say it:} \emph{huāqián}
  \item \emph{Say it:} \emph{huāqián}, once more
  \item \emph{Say it:} \emph{wǒ mèimei hěn huì huāqián} — My little sister really knows how to spend money.
  \item \emph{Recall:} say the Mandarin for in the end, then the Mandarin for well-off, then say \emph{huāqián} again
\end{itemize}

\begin{tcolorbox}[breakable,colback=teal!4,colframe=teal!35!black,title={Before you move on}]
What does huāqián mean? (\textquotedblleft{}To spend money\textquotedblright{}.) Say it once more.
\end{tcolorbox}
\section[língqián]{língqián — small change; loose coins}

\label{lesson:ZH-C188-lingqian}



\begin{quote}
Before the new one: say the Mandarin for well-off, then the Mandarin for to spend money.
\end{quote}

\subsection*{You'll want to know: língqián}
\emph{língqián} — \textquotedblleft{}small change; loose coins\textquotedblright{}.

\begin{quote}
\emph{nǐ yǒu língqián ma} — Have you got any change?
\end{quote}

Its characters come later; for now it is a word for the ear.

\begin{cousinweb}
\emph{língqián} is built from \emph{líng} 'zero' and \zh{钱} 'money', both inherited from Middle Chinese, \emph{leng} and \emph{dzjen} (Guangyun, in Baxter's transcription). The dictionary glosses the first 'to fall' and 'a light rain'; the sense of small odd amounts is not given there, so how the pair came to mean change is not settled.
\end{cousinweb}

\begin{grammarlens}[title={What you've built}]
One more word for asking for what I want.
\end{grammarlens}

\subsection*{Your turn}
\begin{itemize}
\raggedright
  \item \emph{Say it:} \emph{língqián}
  \item \emph{Say it:} \emph{língqián}, once more
  \item \emph{Say it:} \emph{nǐ yǒu língqián ma} — Have you got any change?
  \item \emph{Recall:} say the Mandarin for well-off, then the Mandarin for to spend money, then say \emph{língqián} again
\end{itemize}

\begin{tcolorbox}[breakable,colback=teal!4,colframe=teal!35!black,title={Before you move on}]
What does língqián mean? (\textquotedblleft{}Small change; loose coins\textquotedblright{}.) Say it once more.
\end{tcolorbox}
\section[shěngqián]{shěngqián — to spend less; to save money}

\label{lesson:ZH-C188-shengqian}



\begin{quote}
Before the new one: say the Mandarin for to spend money, then the Mandarin for small change.
\end{quote}

\subsection*{You'll want to know: shěngqián}
\emph{shěngqián} — \textquotedblleft{}to spend less; to save money\textquotedblright{}.

\begin{quote}
\emph{wǒ xiǎng shěngqián suǒyǐ bú qù fàndiàn} — I want to save money, so I'm not going to restaurants.
\end{quote}

Its characters come later; for now it is a word for the ear.

\begin{cousinweb}
\emph{shěngqián} is built from \emph{shěng} and \zh{钱} 'money', both inherited from Middle Chinese, \emph{srjaengX} and \emph{dzjen} (Guangyun, in Baxter's transcription). The dictionary glosses the first as a government office; the sense 'to use less' is not given there, so how it came is not settled.
\end{cousinweb}

\begin{grammarlens}[title={What you've built}]
One more word for asking for what I want.
\end{grammarlens}

\subsection*{Your turn}
\begin{itemize}
\raggedright
  \item \emph{Say it:} \emph{shěngqián}
  \item \emph{Say it:} \emph{shěngqián}, once more
  \item \emph{Say it:} \emph{wǒ xiǎng shěngqián suǒyǐ bú qù fàndiàn} — I want to save money, so I'm not going to restaurants.
  \item \emph{Recall:} say the Mandarin for to spend money, then the Mandarin for small change, then say \emph{shěngqián} again
\end{itemize}

\begin{tcolorbox}[breakable,colback=teal!4,colframe=teal!35!black,title={Before you move on}]
What does shěngqián mean? (\textquotedblleft{}To spend less; to save money\textquotedblright{}.) Say it once more.
\end{tcolorbox}
\section[gòu]{gòu — enough}

\label{lesson:ZH-C188-gou}



\begin{quote}
Before the new one: say the Mandarin for small change, then the Mandarin for to spend less.
\end{quote}

\subsection*{You'll want to know: gòu}
\emph{gòu} — \textquotedblleft{}enough\textquotedblright{}.

\begin{quote}
\emph{wǒ de qián bú gòu} — I haven't got enough money.
\end{quote}

Its characters come later; for now it is a word for the ear.

\begin{cousinweb}
\emph{gòu} is inherited from Middle Chinese \emph{kuw} (Guangyun, in Baxter's transcription), glossed 'many'. The dictionary lists it with a level tone only, so the falling tone of gòu is not the regular outcome, and how it came is not settled.
\end{cousinweb}

\begin{grammarlens}[title={What you've built}]
One more word for asking for what I want.
\end{grammarlens}

\subsection*{Your turn}
\begin{itemize}
\raggedright
  \item \emph{Say it:} \emph{gòu}
  \item \emph{Say it:} \emph{gòu}, once more
  \item \emph{Say it:} \emph{wǒ de qián bú gòu} — I haven't got enough money.
  \item \emph{Recall:} say the Mandarin for small change, then the Mandarin for to spend less, then say \emph{gòu} again
\end{itemize}

\begin{tcolorbox}[breakable,colback=teal!4,colframe=teal!35!black,title={Before you move on}]
What does gòu mean? (\textquotedblleft{}Enough\textquotedblright{}.) Say it once more.
\end{tcolorbox}
